\chapter{Sensibilidad global del diseño v2}\label{cap:sens}

Este capítulo aplica estudios de sensibilidad a las magnitudes que deciden si el diseño
cumple: la velocidad de descenso, el promedio de \Req{C4}, el arranque del rotor, la
autonomía de \Req{E7}, la masa, la altura de liberación, la fuerza centrífuga, el flameo y el
margen frente al anillo de vórtices. El objetivo es saber qué entradas conviene conocer mejor
primero. El resultado principal es que casi todo el riesgo aerodinámico viene de una sola
entrada, el coeficiente de resistencia del rotor $C_R$ (la incertidumbre del modelo), seguida
por la densidad del aire; que el arranque depende sobre todo de la inercia del rotor y del
paso; que \Req{E7} se cumple solo si las cámaras se encienden al despegar; y que el flameo de la
pala depende de dónde queda el centro de masa. Los cálculos están en
\texttt{analysis/v2\_sensibilidad.py}, \texttt{v2\_sens\_arranque.py} y
\texttt{v2\_descenso.py}.

\begin{resultado}[Lo más sensible, en orden]
\begin{enumerate}[leftmargin=*]
  \item $V_2$ y $\langle V\rangle$: el $C_R$ del rotor explica el 79\,\% (Arq.~1) y el 65\,\% (Arq.~2) de
        la varianza de $V_2$; la densidad del aire, el 15 y el 12\,\%; la masa de la etapa 2, el 18\,\%
        en la Arq.~2. Paso, cuerda, número de palas, radio de fabricación y $C_D$ del cuerpo
        suman menos del 8\,\%.
  \item Arranque: la inercia del rotor ($\pm30$\,\% cambia $t_{90}$ de 1,1 a 2,0~s) y el paso
        ($\ang{-2}$ a $\ang{-4}$: 2,0 a 1,3~s). El rozamiento de los rodamientos cambia
        $t_{90}$ en \SI{0.04}{s}.
  \item Autonomía: las cámaras son el \SI{73}{\percent} de la potencia; con todo encendido
        desde la rampa la autonomía es \SI{2.3}{h} y la probabilidad de no llegar a 2\,h es
        del \SI{54}{\percent}.
  \item Masa de la carga útil: el panel solar de la planilla (\SI{37}{g}, dato dudoso) pesa
        el \SI{56}{\percent} de la varianza.
  \item Altura de liberación: el retardo del filtro y el error del puerto de presión
        (\SI{2.3}{m} cada uno); el error de escala de la atmósfera estándar se cancela al usar
        el \SI{80}{\percent} del apogeo medido con el mismo sensor. La altura misma (liberación
        de 400 a \SI{800}{m}) cambia la duración y el peso del arranque, no el promedio de régimen
        (sección~\ref{sec:sens-altura}).
  \item Flameo: la posición del centro de masa de la pala ($\ge0{,}45c$ lo acerca a
        $2{,}5\Omega$); fuerza centrífuga: la velocidad de giro al cuadrado, y el sobrepaso
        del arranque la sube un \SI{60}{\percent}.
  \item Anillo de vórtices: una ráfaga ascendente de \SI{1}{m/s} lleva $V/v_h$ a 1,7.
\end{enumerate}
\end{resultado}

%=====================================================================
\section{Método}\label{sec:sens-metodo}

Para cada respuesta se aplican cuatro herramientas.
\begin{itemize}[leftmargin=*]
  \item \textbf{Elasticidad} local, $E_x=\partial\ln V/\partial\ln x$, por diferencias
        centrales. Dice cuánto cambia la respuesta por cada \SI{1}{\percent} de cambio de $x$.
  \item \textbf{Tornado}: cada entrada recorre su rango plausible con las demás en el valor
        nominal. Ordena el efecto \emph{para ese rango}, que importa tanto como la
        elasticidad.
  \item \textbf{Índices de Sobol} de primer orden $S_i$ y total $S_{T_i}$ (estimadores de
        Saltelli y de Jansen, $2\times10^5$ muestras) para entradas independientes con distribución
        uniforme en el rango: $S_i$ es la parte de la varianza que explica $x_i$ sola, y
        $S_{T_i}-S_i$ la que explica por interacciones. Los $S_i$ suman 1,0 en los casos
        siguientes: el modelo es casi aditivo.
  \item \textbf{PRCC} (correlación parcial de rangos) como control de la monotonía.
\end{itemize}
Las entradas y sus rangos se fijan por ingeniería y se listan en cada sección. La
respuesta de $V_2$ usa el equilibrio $V_2=\sqrt{m_2g/(\tfrac12\rho C_R\pi R^2+\tfrac12\rho C_{D,c}A_c)}$
con $C_R$ del BEMT y sus respuestas lineales al paso, la cuerda y el número de palas medidas en el
barrido del capítulo~\ref{cap:bemt}: $\partial C_R/\partial\theta=\num{0.021}$/\si{\degree},
$\partial C_R/\partial c=\num{0.0017}$/mm y $\partial C_R/\partial B=\num{0.0205}$/pala.

%=====================================================================
\section{Velocidad de régimen $V_2$}\label{sec:sens-V2}

\begin{table}[H]
\centering\small
\caption{Elasticidades y rangos de $V_2$ (Arquitectura~1, $V_2=\SI{6.36}{m/s}$).}\label{tab:sens-elast}
\begin{tabular}{@{}lccc@{}}
\toprule
\textbf{Entrada} & \textbf{Elasticidad} $\partial\ln V_2/\partial\ln x$ & \textbf{Rango} & \textbf{$V_2$ en el rango} [m/s]\\\midrule
$C_R$ del modelo (incertidumbre) & $-0{,}49$ & $\pm15$\,\% & 5,94--6,88\\
Densidad del aire & $-0{,}50$ & 1,10--\SI{1.25}{kg/m^3} & 6,29--6,71\\
Radio de fabricación $R$ & $-0{,}97$ & $\SI{250}{mm}\pm\SI{3}{mm}$ & 6,28--6,43\\
Paso $\theta$ & $-0{,}03$ & \ang{-4} a \ang{-2} & 6,29--6,43\\
Número de palas & $-0{,}04$ & 3 a 5 & 6,29--6,42\\
Masa $m_2$ & $+0{,}50$ & 0,48--\SI{0.50}{kg} & 6,29--6,42\\
Cuerda & $-0{,}04$ & 40 a \SI{50}{mm} & 6,33--6,38\\
$C_D$ del cuerpo & $-0{,}01$ & 0,6--1,0 & 6,34--6,38\\
\bottomrule
\end{tabular}
\end{table}

\begin{figure}[H]
\centering
\begin{tikzpicture}[font=\footnotesize]
\fill[cfijo!70] (0.566,-0.220) rectangle (3.712,0.220);
\fill[crot!80] (3.712,-0.220) rectangle (7.634,0.220);
\node[anchor=east] at (-0.15,-0.000) {modelo: $C_R$ (\textpm\,15\,\%)};
\node[anchor=east, font=\scriptsize] at (0.516,-0.000) {5.94};
\node[anchor=west, font=\scriptsize] at (7.684,-0.000) {6.88};
\fill[cfijo!70] (3.232,-0.840) rectangle (3.712,-0.400);
\fill[crot!80] (3.712,-0.840) rectangle (6.352,-0.400);
\node[anchor=east] at (-0.15,-0.620) {densidad del aire};
\node[anchor=east, font=\scriptsize] at (3.182,-0.620) {6.29};
\node[anchor=west, font=\scriptsize] at (6.402,-0.620) {6.71};
\fill[cfijo!70] (3.160,-1.460) rectangle (3.712,-1.020);
\fill[crot!80] (3.712,-1.460) rectangle (4.276,-1.020);
\node[anchor=east] at (-0.15,-1.240) {radio $R$ (\textpm\,3\,mm)};
\node[anchor=east, font=\scriptsize] at (3.110,-1.240) {6.28};
\node[anchor=west, font=\scriptsize] at (4.326,-1.240) {6.43};
\fill[cfijo!70] (3.213,-2.080) rectangle (3.712,-1.640);
\fill[crot!80] (3.712,-2.080) rectangle (4.227,-1.640);
\node[anchor=east] at (-0.15,-1.860) {paso $\theta$ (\ang{-4} a \ang{-2})};
\node[anchor=east, font=\scriptsize] at (3.163,-1.860) {6.29};
\node[anchor=west, font=\scriptsize] at (4.277,-1.860) {6.43};
\fill[cfijo!70] (3.225,-2.700) rectangle (3.712,-2.260);
\fill[crot!80] (3.712,-2.700) rectangle (4.215,-2.260);
\node[anchor=east] at (-0.15,-2.480) {palas (3 a 5)};
\node[anchor=east, font=\scriptsize] at (3.175,-2.480) {6.29};
\node[anchor=west, font=\scriptsize] at (4.265,-2.480) {6.42};
\fill[cfijo!70] (3.222,-3.320) rectangle (3.712,-2.880);
\fill[crot!80] (3.712,-3.320) rectangle (4.197,-2.880);
\node[anchor=east] at (-0.15,-3.100) {masa de la etapa 2};
\node[anchor=east, font=\scriptsize] at (3.172,-3.100) {6.29};
\node[anchor=west, font=\scriptsize] at (4.247,-3.100) {6.42};
\fill[cfijo!70] (3.508,-3.940) rectangle (3.712,-3.500);
\fill[crot!80] (3.712,-3.940) rectangle (3.918,-3.500);
\node[anchor=east] at (-0.15,-3.720) {cuerda (40 a 50\,mm)};
\node[anchor=east, font=\scriptsize] at (3.458,-3.720) {6.33};
\node[anchor=west, font=\scriptsize] at (3.968,-3.720) {6.38};
\fill[cfijo!70] (3.556,-4.560) rectangle (3.712,-4.120);
\fill[crot!80] (3.712,-4.560) rectangle (3.869,-4.120);
\node[anchor=east] at (-0.15,-4.340) {arrastre de apertura $C_{dc}$};
\node[anchor=east, font=\scriptsize] at (3.506,-4.340) {6.34};
\node[anchor=west, font=\scriptsize] at (3.919,-4.340) {6.38};
\draw[cfreno, dashed, line width=0.6pt] (3.712,0.45) -- (3.712,-4.790);
\node[font=\footnotesize] at (4.100,-5.490) {$V_2$ de régimen [m/s] (Arq.~1)};
\end{tikzpicture}

\caption{Tornado de $V_2$ de régimen, Arquitectura~1 (azul: valor bajo de la entrada;
naranja: valor alto; línea de trazos: nominal).}
\label{fig:sens-tornado-V2}
\end{figure}

La elasticidad del radio es $-0{,}97$: $V_2$ baja casi \SI{1}{\percent} por cada \SI{1}{\percent} de
radio. Pero el radio se fabrica con tolerancia de milímetros, y en el rango de fabricación
pesa \SI{0.15}{m/s}, mientras que la incertidumbre de $C_R$ pesa \SI{0.94}{m/s}. Por eso la
palanca de diseño es el radio y el factor de riesgo es el $C_R$.

\begin{table}[H]
\centering\small
\caption{Índices de Sobol de $V_2$ (primer orden $S_i$; los totales difieren en menos de 0,02).}\label{tab:sens-sobol-V2}
\begin{tabular}{@{}lccc@{}}
\toprule
\textbf{Entrada} & \textbf{Arq.~1} $S_i$ & \textbf{Arq.~2} $S_i$ & PRCC (Arq.~1)\\\midrule
$C_R$ del modelo & 0,79 & 0,65 & $-0{,}99$\\
Densidad del aire & 0,15 & 0,12 & $-0{,}77$\\
Masa $m_2$ & 0,01 & 0,18 & $+0{,}33$\\
Radio, paso, cuerda, palas, $C_{D,c}$ & 0,05 & 0,05 & ---\\
\bottomrule
\end{tabular}
\end{table}

En la Arquitectura~2 la masa pesa más porque su incertidumbre es mayor (\SI{15}{g}, de la
suma de componentes, contra \SI{6}{g} del pesaje en la Arquitectura~1). \textbf{Consecuencia
para los ensayos:} medir $C_R$ (caída instrumentada, ensayo E3 del capítulo~\ref{cap:ensayos}) reduce la
varianza de $V_2$ en tres cuartas partes; medir el paso o el radio con más precisión casi no
aporta.

\begin{figure}[H]
\centering
\begin{tikzpicture}[font=\footnotesize]
\fill[cfijo!70] (0.566,-0.220) rectangle (3.712,0.220);
\fill[crot!80] (3.712,-0.220) rectangle (7.634,0.220);
\node[anchor=east] at (-0.15,-0.000) {modelo: $C_R$ (\textpm\,15\,\%)};
\node[anchor=east, font=\scriptsize] at (0.516,-0.000) {4.84};
\node[anchor=west, font=\scriptsize] at (7.684,-0.000) {5.60};
\fill[cfijo!70] (1.839,-0.840) rectangle (3.712,-0.400);
\fill[crot!80] (3.712,-0.840) rectangle (5.514,-0.400);
\node[anchor=east] at (-0.15,-0.620) {masa de la etapa 2};
\node[anchor=east, font=\scriptsize] at (1.789,-0.620) {4.97};
\node[anchor=west, font=\scriptsize] at (5.564,-0.620) {5.37};
\fill[cfijo!70] (3.232,-1.460) rectangle (3.712,-1.020);
\fill[crot!80] (3.712,-1.460) rectangle (6.352,-1.020);
\node[anchor=east] at (-0.15,-1.240) {densidad del aire};
\node[anchor=east, font=\scriptsize] at (3.182,-1.240) {5.13};
\node[anchor=west, font=\scriptsize] at (6.402,-1.240) {5.46};
\fill[cfijo!70] (3.160,-2.080) rectangle (3.712,-1.640);
\fill[crot!80] (3.712,-2.080) rectangle (4.276,-1.640);
\node[anchor=east] at (-0.15,-1.860) {radio $R$ (\textpm\,3\,mm)};
\node[anchor=east, font=\scriptsize] at (3.110,-1.860) {5.12};
\node[anchor=west, font=\scriptsize] at (4.326,-1.860) {5.24};
\fill[cfijo!70] (3.213,-2.700) rectangle (3.712,-2.260);
\fill[crot!80] (3.712,-2.700) rectangle (4.227,-2.260);
\node[anchor=east] at (-0.15,-2.480) {paso $\theta$ (\ang{-4} a \ang{-2})};
\node[anchor=east, font=\scriptsize] at (3.163,-2.480) {5.12};
\node[anchor=west, font=\scriptsize] at (4.277,-2.480) {5.23};
\fill[cfijo!70] (3.225,-3.320) rectangle (3.712,-2.880);
\fill[crot!80] (3.712,-3.320) rectangle (4.215,-2.880);
\node[anchor=east] at (-0.15,-3.100) {palas (3 a 5)};
\node[anchor=east, font=\scriptsize] at (3.175,-3.100) {5.12};
\node[anchor=west, font=\scriptsize] at (4.265,-3.100) {5.23};
\fill[cfijo!70] (3.508,-3.940) rectangle (3.712,-3.500);
\fill[crot!80] (3.712,-3.940) rectangle (3.919,-3.500);
\node[anchor=east] at (-0.15,-3.720) {cuerda (40 a 50\,mm)};
\node[anchor=east, font=\scriptsize] at (3.458,-3.720) {5.16};
\node[anchor=west, font=\scriptsize] at (3.969,-3.720) {5.20};
\fill[cfijo!70] (3.556,-4.560) rectangle (3.712,-4.120);
\fill[crot!80] (3.712,-4.560) rectangle (3.869,-4.120);
\node[anchor=east] at (-0.15,-4.340) {arrastre de apertura $C_{dc}$};
\node[anchor=east, font=\scriptsize] at (3.506,-4.340) {5.16};
\node[anchor=west, font=\scriptsize] at (3.919,-4.340) {5.19};
\draw[cfreno, dashed, line width=0.6pt] (3.712,0.45) -- (3.712,-4.790);
\node[font=\footnotesize] at (4.100,-5.490) {$V_2$ de régimen [m/s] (Arq.~2)};
\end{tikzpicture}

\caption{Tornado de $V_2$, Arquitectura~2 ($m_2=\SI{0.325}{kg}$).}
\label{fig:sens-tornado-V2b}
\end{figure}

%=====================================================================
\section{Promedio de \texorpdfstring{\Req{C4}}{C4} sobre todo el descenso}\label{sec:sens-Vg}

El promedio global~\eqref{eq:promedio}, con el apogeo de diseño (sección~\ref{sec:dim-tiempos}),
suma a lo anterior las entradas de la etapa 1 (masa de pesaje, $C_D$ y diámetro del drogue), la
fracción de altura recorrida a $V_1$ (el error del disparo del 80\,\%, de 17 a 23\,\%) y el
tiempo de transición de $V_1$ a $V_2$ (2,5 a \SI{5}{s}; sección~\ref{sec:arq-mc}).

\begin{figure}[H]
\centering
\begin{tikzpicture}[font=\footnotesize]
\fill[cfijo!70] (0.566,-0.220) rectangle (3.745,0.220);
\fill[crot!80] (3.745,-0.220) rectangle (7.634,0.220);
\node[anchor=east] at (-0.15,-0.000) {modelo: $C_R$ (\textpm\,15\,\%)};
\node[anchor=east, font=\scriptsize] at (0.516,-0.000) {6.95};
\node[anchor=west, font=\scriptsize] at (7.684,-0.000) {7.91};
\fill[cfijo!70] (2.046,-0.840) rectangle (3.745,-0.400);
\fill[crot!80] (3.745,-0.840) rectangle (5.616,-0.400);
\node[anchor=east] at (-0.15,-0.620) {densidad del aire};
\node[anchor=east, font=\scriptsize] at (1.996,-0.620) {7.15};
\node[anchor=west, font=\scriptsize] at (5.666,-0.620) {7.63};
\fill[cfijo!70] (2.789,-1.460) rectangle (3.745,-1.020);
\fill[crot!80] (3.745,-1.460) rectangle (4.734,-1.020);
\node[anchor=east] at (-0.15,-1.240) {fracción a $V_1$ (disparo)};
\node[anchor=east, font=\scriptsize] at (2.739,-1.240) {7.25};
\node[anchor=west, font=\scriptsize] at (4.784,-1.240) {7.51};
\fill[cfijo!70] (3.244,-2.080) rectangle (3.745,-1.640);
\fill[crot!80] (3.745,-2.080) rectangle (4.260,-1.640);
\node[anchor=east] at (-0.15,-1.860) {paso $\theta$ (\ang{-4} a \ang{-2})};
\node[anchor=east, font=\scriptsize] at (3.194,-1.860) {7.31};
\node[anchor=west, font=\scriptsize] at (4.310,-1.860) {7.45};
\fill[cfijo!70] (3.253,-2.700) rectangle (3.745,-2.260);
\fill[crot!80] (3.745,-2.700) rectangle (4.230,-2.260);
\node[anchor=east] at (-0.15,-2.480) {masa de la etapa 2};
\node[anchor=east, font=\scriptsize] at (3.203,-2.480) {7.32};
\node[anchor=west, font=\scriptsize] at (4.280,-2.480) {7.45};
\fill[cfijo!70] (3.404,-3.320) rectangle (3.745,-2.880);
\fill[crot!80] (3.745,-3.320) rectangle (4.089,-2.880);
\node[anchor=east] at (-0.15,-3.100) {tiempo de arranque};
\node[anchor=east, font=\scriptsize] at (3.354,-3.100) {7.34};
\node[anchor=west, font=\scriptsize] at (4.139,-3.100) {7.43};
\fill[cfijo!70] (3.435,-3.940) rectangle (3.745,-3.500);
\fill[crot!80] (3.745,-3.940) rectangle (4.078,-3.500);
\node[anchor=east] at (-0.15,-3.720) {$C_D$ del drogue};
\node[anchor=east, font=\scriptsize] at (3.385,-3.720) {7.34};
\node[anchor=west, font=\scriptsize] at (4.128,-3.720) {7.43};
\fill[cfijo!70] (3.465,-4.560) rectangle (3.745,-4.120);
\fill[crot!80] (3.745,-4.560) rectangle (4.030,-4.120);
\node[anchor=east] at (-0.15,-4.340) {diámetro del drogue};
\node[anchor=east, font=\scriptsize] at (3.415,-4.340) {7.34};
\node[anchor=west, font=\scriptsize] at (4.080,-4.340) {7.42};
\fill[cfijo!70] (3.480,-5.180) rectangle (3.745,-4.740);
\fill[crot!80] (3.745,-5.180) rectangle (4.015,-4.740);
\node[anchor=east] at (-0.15,-4.960) {arrastre de apertura $C_{dc}$};
\node[anchor=east, font=\scriptsize] at (3.430,-4.960) {7.35};
\node[anchor=west, font=\scriptsize] at (4.065,-4.960) {7.42};
\draw[cfreno, dashed, line width=0.6pt] (3.745,0.45) -- (3.745,-5.410);
\draw[cgris, dotted, line width=0.8pt] (8.335,0.45) -- (8.335,-5.410) node[below, font=\scriptsize, cgris] {límite 8};
\node[font=\footnotesize] at (4.100,-6.110) {velocidad media global [m/s] (Arq.~1)};
\end{tikzpicture}

\caption{Tornado del promedio global, Arquitectura~1 (nominal \SI{7.42}{m/s}; el límite de
\SI{8}{m/s} se indica con puntos).}
\label{fig:sens-tornado-Vg}
\end{figure}

\begin{table}[H]
\centering\small
\caption{Índices de Sobol del promedio global.}\label{tab:sens-sobol-Vg}
\begin{tabular}{@{}lcc@{}}
\toprule
\textbf{Entrada} & \textbf{Arq.~1} $S_i$ & \textbf{Arq.~2} $S_i$\\\midrule
$C_R$ del modelo & 0,70 & 0,59\\
Densidad del aire & 0,18 & 0,15\\
Masa $m_2$ & 0,01 & 0,16\\
Fracción de altura a $V_1$ (disparo del 80\,\%) & 0,05 & 0,06\\
Etapa 1 ($m_1$, $C_{D,d}$, $D_d$, $C_{D,c}$) y transición & 0,03 & 0,02\\
\bottomrule
\end{tabular}
\end{table}

La etapa 1 casi no importa para el promedio: $m_1$, $C_D$ y el diámetro del drogue suman menos
del \SI{2}{\percent} de la varianza. La probabilidad de superar \SI{8}{m/s} con estas entradas
es del \SI{4.8}{\percent} para la Arquitectura~1 y nula para la 2, menor que el
\SI{14}{\percent} del Monte Carlo de la sección~\ref{sec:arq-mc} porque allí la densidad se
muestrea con la altura del sitio y la temperatura (más cola superior); ambos usan el mismo
apogeo.

%=====================================================================
\section{Arranque del rotor}\label{sec:sens-arranque}

\begin{figure}[H]
\centering
\begin{tikzpicture}[font=\footnotesize]
\fill[cfijo!70] (0.566,-0.220) rectangle (4.031,0.220);
\fill[crot!80] (4.031,-0.220) rectangle (7.542,0.220);
\node[anchor=east] at (-0.15,-0.000) {inercia del rotor (\textpm\,30\,\%)};
\node[anchor=east, font=\scriptsize] at (0.516,-0.000) {1.10};
\node[anchor=west, font=\scriptsize] at (7.592,-0.000) {2.01};
\fill[cfijo!70] (2.007,-0.840) rectangle (4.031,-0.400);
\fill[crot!80] (4.031,-0.840) rectangle (7.634,-0.400);
\node[anchor=east] at (-0.15,-0.620) {paso $\theta$ (\ang{-4} a \ang{-2})};
\node[anchor=east, font=\scriptsize] at (1.957,-0.620) {1.29};
\node[anchor=west, font=\scriptsize] at (7.684,-0.620) {2.02};
\fill[cfijo!70] (2.758,-1.460) rectangle (4.031,-1.020);
\fill[crot!80] (4.031,-1.460) rectangle (5.948,-1.020);
\node[anchor=east] at (-0.15,-1.240) {$V$ al abrir (10 a 16\,m/s)};
\node[anchor=east, font=\scriptsize] at (2.708,-1.240) {1.39};
\node[anchor=west, font=\scriptsize] at (5.998,-1.240) {1.80};
\fill[cfijo!70] (2.682,-2.080) rectangle (4.031,-1.640);
\fill[crot!80] (4.031,-2.080) rectangle (4.031,-1.640);
\node[anchor=east] at (-0.15,-1.860) {polar en pérdida (XFOIL o 2D)};
\node[anchor=east, font=\scriptsize] at (2.632,-1.860) {1.38};
\node[anchor=west, font=\scriptsize] at (4.081,-1.860) {1.55};
\fill[cfijo!70] (4.031,-2.700) rectangle (4.031,-2.260);
\fill[crot!80] (4.031,-2.700) rectangle (4.782,-2.260);
\node[anchor=east] at (-0.15,-2.480) {masa de la etapa 2};
\node[anchor=east, font=\scriptsize] at (3.981,-2.480) {1.55};
\node[anchor=west, font=\scriptsize] at (4.832,-2.480) {1.65};
\fill[cfijo!70] (3.954,-3.320) rectangle (4.031,-2.880);
\fill[crot!80] (4.031,-3.320) rectangle (4.230,-2.880);
\node[anchor=east] at (-0.15,-3.100) {rozamiento $Q_f$ (0,34 a 1,7\,mN.m)};
\node[anchor=east, font=\scriptsize] at (3.904,-3.100) {1.54};
\node[anchor=west, font=\scriptsize] at (4.280,-3.100) {1.58};
\fill[cfijo!70] (4.000,-3.940) rectangle (4.031,-3.500);
\fill[crot!80] (4.031,-3.940) rectangle (4.031,-3.500);
\node[anchor=east] at (-0.15,-3.720) {arrastre de apertura $C_{dc}$};
\node[anchor=east, font=\scriptsize] at (3.950,-3.720) {1.55};
\node[anchor=west, font=\scriptsize] at (4.081,-3.720) {1.55};
\draw[cfreno, dashed, line width=0.6pt] (4.031,0.45) -- (4.031,-4.170);
\node[font=\footnotesize] at (4.100,-4.870) {$t_{90}$ del arranque [s]};
\end{tikzpicture}

\caption{Tornado del tiempo de arranque $t_{90}$ (capítulo~\ref{cap:arranque}, caso base
\SI{1.55}{s}).}
\label{fig:sens-tornado-t90}
\end{figure}

La inercia del rotor es la entrada más sensible: un \SI{30}{\percent} más de inercia
alarga el arranque de 1,55 a \SI{2.0}{s} y la caída de 22 a \SI{28}{m} (del \SI{4.3}{\percent} al
\SI{5.6}{\percent} de la altura de la etapa~2). Como la inercia
depende de la masa de las palas (portapala y varilla de BV incluidos) y de su centro de masa,
\textbf{pesar y balancear las palas} es la medida más barata. Le sigue el paso: de \ang{-2} a
\ang{-4} el arranque va de 2,0 a \SI{1.3}{s}. La polar en pérdida (XFOIL contra 2D) cambia
\SI{0.17}{s}; la velocidad al abrir, \SI{0.4}{s}; el rozamiento, \SI{0.04}{s}; la masa de la
etapa 2 y el arrastre de la apertura, menos de \SI{0.1}{s}. En cambio la \emph{existencia} del
arranque (paso crítico $\theta_c\approx\ang{0}$) casi no depende de ninguno: solo importa que
el paso sea negativo.

%=====================================================================
\section{Autonomía y \texorpdfstring{\Req{E7}}{E7}}\label{sec:sens-energia}

\begin{figure}[H]
\centering
\begin{tikzpicture}[font=\footnotesize]
\fill[cfijo!70] (1.517,-0.220) rectangle (2.368,0.220);
\fill[crot!80] (2.368,-0.220) rectangle (7.634,0.220);
\node[anchor=east] at (-0.15,-0.000) {cámaras (c/u)};
\node[anchor=east, font=\scriptsize] at (1.467,-0.000) {2.21};
\node[anchor=west, font=\scriptsize] at (7.684,-0.000) {3.42};
\fill[cfijo!70] (1.033,-0.840) rectangle (2.368,-0.400);
\fill[crot!80] (2.368,-0.840) rectangle (3.035,-0.400);
\node[anchor=east] at (-0.15,-0.620) {fracción útil de la celda};
\node[anchor=east, font=\scriptsize] at (0.983,-0.620) {2.11};
\node[anchor=west, font=\scriptsize] at (3.085,-0.620) {2.51};
\fill[cfijo!70] (0.566,-1.460) rectangle (2.368,-1.020);
\fill[crot!80] (2.368,-1.460) rectangle (2.368,-1.020);
\node[anchor=east] at (-0.15,-1.240) {efecto de la temperatura};
\node[anchor=east, font=\scriptsize] at (0.516,-1.240) {2.02};
\node[anchor=west, font=\scriptsize] at (2.418,-1.240) {2.38};
\fill[cfijo!70] (0.601,-2.080) rectangle (2.368,-1.640);
\fill[crot!80] (2.368,-2.080) rectangle (2.368,-1.640);
\node[anchor=east] at (-0.15,-1.860) {capacidad de la celda};
\node[anchor=east, font=\scriptsize] at (0.551,-1.860) {2.03};
\node[anchor=west, font=\scriptsize] at (2.418,-1.860) {2.38};
\fill[cfijo!70] (1.226,-2.700) rectangle (2.368,-2.260);
\fill[crot!80] (2.368,-2.700) rectangle (2.917,-2.260);
\node[anchor=east] at (-0.15,-2.480) {rendimiento elevador};
\node[anchor=east, font=\scriptsize] at (1.176,-2.480) {2.15};
\node[anchor=west, font=\scriptsize] at (2.967,-2.480) {2.49};
\fill[cfijo!70] (2.044,-3.320) rectangle (2.368,-2.880);
\fill[crot!80] (2.368,-3.320) rectangle (3.234,-2.880);
\node[anchor=east] at (-0.15,-3.100) {radio (LoRa)};
\node[anchor=east, font=\scriptsize] at (1.994,-3.100) {2.31};
\node[anchor=west, font=\scriptsize] at (3.284,-3.100) {2.55};
\fill[cfijo!70] (1.365,-3.940) rectangle (2.368,-3.500);
\fill[crot!80] (2.368,-3.940) rectangle (2.512,-3.500);
\node[anchor=east] at (-0.15,-3.720) {MCU};
\node[anchor=east, font=\scriptsize] at (1.315,-3.720) {2.18};
\node[anchor=west, font=\scriptsize] at (2.562,-3.720) {2.41};
\fill[cfijo!70] (2.089,-4.560) rectangle (2.368,-4.120);
\fill[crot!80] (2.368,-4.560) rectangle (2.561,-4.120);
\node[anchor=east] at (-0.15,-4.340) {microSD};
\node[anchor=east, font=\scriptsize] at (2.039,-4.340) {2.32};
\node[anchor=west, font=\scriptsize] at (2.611,-4.340) {2.42};
\fill[cfijo!70] (2.070,-5.180) rectangle (2.368,-4.740);
\fill[crot!80] (2.368,-5.180) rectangle (2.522,-4.740);
\node[anchor=east] at (-0.15,-4.960) {servos en reposo};
\node[anchor=east, font=\scriptsize] at (2.020,-4.960) {2.32};
\node[anchor=west, font=\scriptsize] at (2.572,-4.960) {2.41};
\draw[cfreno, dashed, line width=0.6pt] (2.368,0.45) -- (2.368,-5.410);
\draw[cgris, dotted, line width=0.8pt] (0.458,0.45) -- (0.458,-5.410) node[below, font=\scriptsize, cgris] {E7: 2 h};
\node[font=\footnotesize] at (4.100,-6.110) {autonomía con todo encendido [h]};
\end{tikzpicture}

\caption{Tornado de la autonomía con todo encendido desde la rampa (planilla del equipo;
topología batería $\to$ elevador a 5\,V $\to$ reductor a 3,3\,V; nominal \SI{2.29}{h}).}
\label{fig:sens-tornado-E7}
\end{figure}

El consumo de las cámaras (\SI{300}{mA} a \SI{5}{V} cada una, \SI{3}{W} en total) es el
\SI{73}{\percent} de la potencia de la batería. Con todo encendido desde la rampa y los
valores de la planilla, la autonomía es de \SI{2.29}{h} y el Monte Carlo de las trece
entradas da \SI{2.0}{h} de media, \SI{1.5}{h} en el percentil 5 y un \SI{54}{\percent} de
probabilidad de no llegar a las 2\,h de \Req{E7} (\SI{53}{\percent} con el Monte Carlo del capítulo~\ref{cap:sensores}, que usa otra muestra y otras distribuciones de entrada: son la misma estimación). Cambiar de \SI{300}{mA} a \SI{180}{mA} por
cámara sube la autonomía a \SI{3.2}{h}; las pérdidas de la doble conversión (0,77) cuestan
unos \SI{0.1}{h}.

\begin{table}[H]
\centering\small
\caption{Autonomía según la política de energía (celda de \SI{3400}{mAh}, \SI{90}{\percent} útil).}\label{tab:sens-energia}
\begin{tabular}{@{}lccc@{}}
\toprule
\textbf{Política} & \textbf{Potencia} [W] & \textbf{Autonomía} [h] & \Req{E7}\\\midrule
Todo encendido desde la rampa, 3,3\,V vía 5\,V & 4,81 & 2,3 & Justo\\
Ídem, 3,3\,V directo de la batería (buck-boost) & 4,63 & 2,4 & Justo\\
Cámaras apagadas en la rampa, radio continua (105\,mA) & 1,28 & 8,6 & Cumple\\
Cámaras apagadas en la rampa, radio a 1\,Hz (20\,mA) & 0,91 & 12,1 & Cumple\\
\bottomrule
\end{tabular}
\end{table}

La conclusión de diseño es encender las cámaras al detectar el despegue (por aceleración o por
presión) y mantener la rampa a $\approx\SI{1}{W}$. Con esa política las 2 horas de \Req{E7} se cubren con un factor
4 (rampa) más el vuelo con las cámaras encendidas: unos \SIrange{88}{101}{s} de descenso
(sección~\ref{sec:dim-tiempos}) y la subida, de decenas de segundos.

%=====================================================================
\section{Masa, altura de liberación y fuerza centrífuga}\label{sec:sens-masa}

\paragraph{Masa.} La tabla~\ref{tab:masa} del capítulo~\ref{cap:dim} suma $497\pm21$~g
(Arq.~1) antes del lastre, con los cambios de electrónica, la cabeza del rotor corregida y las aletas
(sección~\ref{sec:arq-integracion}). Si todos los componentes se pesan con las desviaciones de la
tabla, el lastre necesario es de $3\pm21$~g y hay \SI{44}{\percent} de probabilidad de
sobrepeso (lastre negativo); con un panel de \SI{8}{g} el lastre es de $32\pm14$~g y el sobrepeso, de
\SI{1.9}{\percent}. $V_2$ queda
en $6{,}36\pm0{,}04$~m/s. El panel solar explica la mayor parte de la varianza (planilla:
\SI{37}{g}; típico: \SI{6}{g}), y la segunda celda de \SI{48}{g} no es viable.

\paragraph{Altura de liberación.} El error de la altura a la que se suelta el rotor
resulta de: retardo del filtro $\times$ $V_1$ (\SI{2.3}{m}), error del puerto estático por
la presión dinámica del aire a \SI{13}{m/s} (\SI{2.4}{m}), deriva térmica del offset del
BMP280 (\SI{0.7}{m}), deriva meteorológica en la rampa (\SI{0.4}{m}), ruido (\SI{0.15}{m}) y
pico de ruido al buscar el máximo (\SI{0.18}{m}). Suman un error medio de $+\SI{7.2}{m}$
(se libera más abajo) con \SI{3.4}{m} de desviación y \SI{12.9}{m} en el percentil 95: el
\SI{2.6}{\percent} de la altura de liberación y el \SI{2.1}{\percent} del apogeo. Eso mueve el
promedio global menos de \SI{0.09}{m/s}. Si el umbral fuera absoluto, el error de escala de la
atmósfera estándar ($\pm4$\,\%) daría \SI{20}{m} a la altura de liberación; al usar el
\SI{80}{\percent} del apogeo medido con el mismo sensor, ese error se cancela. El diseño del puerto de presión (venteado, lejos de la corriente) y la
frecuencia de muestreo importan más que la calidad del sensor.

\paragraph{Fuerza centrífuga y flameo.} La fuerza centrífuga por pala vale
$F_c=m_b\Omega^2r_g$. La velocidad de giro es la entrada que domina: el sobrepaso del
arranque ($\Omega=\SI{162.6}{rad/s}$ contra \SI{128.7}{rad/s} de régimen) multiplica $F_c$ por
1,60 (de \SI{18.9}{N} a \SI{30.2}{N} en la pala A). Una incertidumbre de \SI{15}{\percent} en
$m_b$ cambia $F_c$ lo mismo. Para el flameo, la posición del centro de masa de la pala domina:

\begin{table}[H]
\centering\small
\caption{Velocidad de flameo $\Omega_F/\Omega$ (a \SI{128.7}{rad/s}) según la posición del
centro de masa $y_{cm}/c$ (modelo p-k del capítulo~\ref{cap:estructura}).}\label{tab:sens-flameo}
\begin{tabular}{@{}lccccc@{}}
\toprule
$y_{cm}/c$ & 0,40 & 0,425 & 0,45 & 0,50 & 0,55\\\midrule
Pala A & $\infty$ & 1,47 & 1,17 & 0,97 & 0,86\\
Pala BV & 3,34 & 2,75 & 2,47 & 2,20 & 2,12\\
\bottomrule
\end{tabular}
\end{table}

Con BV, un centro de masa corrido de $0{,}44c$ a $0{,}55c$ baja el margen de 2,6 a 2,1; con A,
cualquier posición por detrás de $0{,}41c$ da flameo en régimen o muy cerca. El centro de masa
se mide con el ensayo de balanceo sobre filo (capítulo~\ref{cap:estructura}).

%=====================================================================
\section{Altura de liberación y apogeo}\label{sec:sens-altura}

El equipo estima la liberación del rotor a \SI{500}{m}; con el \SI{80}{\percent} de \Req{C4}, el
apogeo de diseño es de \SI{625}{m} (caso base, sección~\ref{sec:dim-tiempos}). Si lo estimado
fuera el apogeo, la liberación sería a \SI{400}{m} (lectura alternativa); \SI{800}{m} (apogeo de
\SI{1000}{m}) es una cota superior. El promedio de régimen no cambia (\SI{7.11}{m/s} con la
Arq.~1 y \SI{6.21}{m/s} con la 2). Cambian la duración, el peso del arranque y el error relativo
de la liberación.

\begin{table}[H]
\centering\small
\caption{Sensibilidad a la altura de liberación (caso base en la columna central). El promedio con
transición supone $t_s=\SI{3.75}{s}$; las probabilidades salen del Monte Carlo de la
sección~\ref{sec:arq-mc} con el apogeo fijo y los mismos sorteos en las tres alturas.}\label{tab:sens-altura}
\begin{tabular}{@{}lccc@{}}
\toprule
Altura de liberación [m] & 400 & \textbf{500} & 800\\
Apogeo [m] & 500 & \textbf{625} & 1000\\\midrule
Etapa 1 [s] & 7,4 & \textbf{9,3} & 14,9\\
Etapa 2, Arq.~1 [s] & 62,9 & \textbf{78,6} & 125,8\\
Etapa 2, Arq.~2 [s] & 73,1 & \textbf{91,4} & 146,3\\
$\langle V\rangle$ con transición, Arq.~1 [m/s] & 7,33 & \textbf{7,28} & 7,22\\
$\langle V\rangle$ con transición, Arq.~2 [m/s] & 6,42 & \textbf{6,38} & 6,31\\
$P(\text{C4, global})$, Arq.~1 [\%] & 83,8 & \textbf{86,0} & 88,9\\
$P(\text{C4, global})$, Arq.~2 [\%] & 99,8 & \textbf{99,8} & 99,9\\
Error de liberación p95 (\SI{12.9}{m}) [\% de la altura de liberación] & 3,2 & \textbf{2,6} & 1,6\\
Caída del arranque (\SI{21.7}{m}) [\% de la altura de liberación] & 5,4 & \textbf{4,3} & 2,7\\
\bottomrule
\end{tabular}
\end{table}

Una liberación más baja alarga proporcionalmente el arranque y el error de la liberación, y
sube el promedio con transición (la transición se recorre más rápido que $V_2$ y pesa más): la
probabilidad de cumplir baja de \SI{88.9}{\percent} a \SI{83.8}{\percent} en la Arq.~1 entre
\SI{800}{m} y \SI{400}{m}, y casi no cambia en la 2. La tabla no incluye la menor densidad a la
altura de vuelo (sección~\ref{sec:arq-mc}), que baja esas probabilidades y crece con la altura.

%=====================================================================
\section{Anillo de vórtices y ráfagas}\label{sec:sens-vrs}

El cociente $V/v_h$ en régimen es 2,04, que no depende de la masa ni del radio: la
sensibilidad es solo a la ráfaga vertical $w$.
\begin{table}[H]
\centering\small
\caption{Ráfaga ascendente $w$ necesaria para llevar $V/v_h$ a 1,7 y a 1,5.}\label{tab:sens-gust}
\begin{tabular}{@{}lccc@{}}
\toprule
\textbf{Arquitectura} & $v_h$ [m/s] & $w$ hasta $1{,}7\,v_h$ [m/s] & $w$ hasta $1{,}5\,v_h$ [m/s]\\\midrule
1 ($m_2=\SI{0.49}{kg}$) & 3,12 & 1,05 & 1,68\\
2 ($m_2=\SI{0.325}{kg}$) & 2,54 & 0,86 & 1,37\\
\bottomrule
\end{tabular}
\end{table}
La Arquitectura~2 tiene menos margen absoluto porque $v_h$ es menor. Ambos valores son
pocos para turbulencia fuerte, pero la perturbación dura la constante de tiempo vertical
(\SI{0.3}{s}) y el sistema se aleja solo de la región (capítulo~\ref{cap:estabilidad}).

%=====================================================================
\section{Qué medir primero}\label{sec:sens-prioridad}

\begin{table}[H]
\centering\small
\caption{Prioridad de mediciones según la sensibilidad.}\label{tab:sens-prioridad}
\begin{tabularx}{\textwidth}{@{}lYY@{}}
\toprule
\textbf{Respuesta} & \textbf{Entrada dominante} & \textbf{Qué hacer}\\\midrule
$V_2$, $\langle V\rangle$ & $C_R$ (60--80\,\% de la varianza); luego densidad & Caída instrumentada con dron (E3); fijar la altura del sitio\\
Arranque, $t_{90}$ & Inercia del rotor; paso & Pesar y balancear las palas; portapalas de \ang{-3} a \ang{-4}\\
Autonomía & Cámaras (73\,\% de la potencia) & Encenderlas al despegar; medir su consumo real\\
Masa & Panel solar (56\,\%) & Pesarlo antes que nada\\
Altura de liberación & Retardo del filtro y puerto estático & Ventear el puerto; muestrear a \SI{20}{Hz} o más\\
Flameo & Posición del centro de masa & Balanceo sobre filo de cada pala\\
Fuerza centrífuga & Sobrepaso de $\Omega$ & Diseñar a \num{1553}~rpm\\
Anillo de vórtices & Ráfaga & Sin acción; vigilar $V/v_h$ en los ensayos\\
\bottomrule
\end{tabularx}
\end{table}
